% =====================================================================
% Test d'espacement vertical entre boites — lu par check-spacing.sh
%
% Chaque cas mesure, avec \pdfsavepos, l'ecart vertical entre la derniere
% ligne d'une boite et la premiere ligne de la boite qui la suit, et l'ecrit
% dans le journal sous la forme « OCOTS-GAP <cas> <ecart en sp> ». Le cas
% `ref' (rien entre les deux boites) sert d'etalon : tout autre cas doit
% donner le meme ecart. Le defaut est silencieux (pas d'erreur LaTeX), d'ou
% ce test.
%
% Un cas par page, en \raggedbottom : une note de bas de page ou l'etirement
% d'une page justifiee ne doivent pas deplacer ce qu'on mesure.
% =====================================================================
\documentclass[11pt]{ocots-book}
\usepackage[lang=fr]{ocots}
\raggedbottom

\makeatletter
% Position verticale de la ligne courante, relue a l'expedition de la page.
\newcommand{\gapmark}[1]{\pdfsavepos\write-1{OCOTS-POS #1 \the\pdflastypos}}
\makeatother

\begin{document}
\mainmatter
\chapter{Espacement}

% ---------------------------------------------------------------------
% ref : boite puis preuve, rien entre les deux.
\begin{proposition}
    Enonce.\gapmark{ref-a}
\end{proposition}
\begin{proof}
    \gapmark{ref-b}Preuve.
\end{proof}
\clearpage

% ---------------------------------------------------------------------
% footnotetext : un \footnotetext entre la boite et la preuve
% (ocots-latex-template#51).
\begin{proposition}
    Enonce\footnotemark{}.\gapmark{footnotetext-a}
\end{proposition}
\footnotetext{La note de l'enonce.}
\begin{proof}
    \gapmark{footnotetext-b}Preuve.
\end{proof}
\clearpage

% ---------------------------------------------------------------------
% footnotetext-box : meme chose entre deux boites de resultat.
\begin{proposition}
    Enonce\footnotemark{}.\gapmark{footnotetext-box-a}
\end{proposition}
\footnotetext{La note de l'enonce.}
\begin{proposition}
    \gapmark{footnotetext-box-b}Enonce.
\end{proposition}
\clearpage

% ---------------------------------------------------------------------
% ref-box : etalon pour deux boites de resultat consecutives.
\begin{proposition}
    Enonce.\gapmark{ref-box-a}
\end{proposition}
\begin{proposition}
    \gapmark{ref-box-b}Enonce.
\end{proposition}

\end{document}
